\documentclass[
	% -- opções da classe memoir --
	12pt,				% tamanho da fonte
	openright,			% capítulos começam em pág ímpar (insere página vazia caso preciso)
	oneside,			% para impressão em verso e anverso. Oposto a oneside
	a4paper,			% tamanho do papel. 
	% -- opções da classe abntex2 --
	chapter=TITLE,		% títulos de capítulos convertidos em letras maiúsculas
	%section=TITLE,		% títulos de seções convertidos em letras maiúsculas
	subsection=TITLE,	% títulos de subseções convertidos em letras maiúsculas
	%subsubsection=TITLE,% títulos de subsubseções convertidos em letras maiúsculas
	% -- opções do pacote babel --
	english,			% idioma adicional para hifenização
	brazil,				% o último idioma é o principal do documento
	]{abntex2}

% ---
% PACOTES
% ---

% ---
% Pacotes fundamentais 
% ---
\usepackage{lmodern}			% Usa a fonte Latin Modern
\usepackage[T1]{fontenc}		% Selecao de codigos de fonte.
\usepackage[utf8]{inputenc}		% Codificacao do documento (conversão automática dos acentos)
\usepackage{indentfirst}		% Indenta o primeiro parágrafo de cada seção.
\usepackage{color}				% Controle das cores
\usepackage{graphicx}			% Inclusão de gráficos
\usepackage{microtype} 			% para melhorias de justificação

% ---

%--------------Pacotes de citações
\usepackage[brazilian,hyperpageref]{backref}	 % Paginas com as citações na bibl
\usepackage[alf]{abntex2cite}	% Citações padrão ABNT

\usepackage{hyperref}

% --- 
% CONFIGURAÇÕES DE PACOTES
% --- 

% ---
% Configurações do pacote backref
% Usado sem a opção hyperpageref de backref
\renewcommand{\backrefpagesname}{Citado na(s) página(s):~}
% Texto padrão antes do número das páginas
\renewcommand{\backref}{}
% Define os textos da citação
\renewcommand*{\backrefalt}[4]{
	\ifcase #1 %
		Nenhuma citação no texto.%
	\or
		Citado na página #2.%
	\else
		Citado #1 vezes nas páginas #2.%
	\fi}%
% ---

% alterando o aspecto da cor azul
\definecolor{blue}{RGB}{41,5,195}

% informações do PDF
\makeatletter
\hypersetup{
     	%pagebackref=true,
		pdftitle={\@title}, 
		pdfauthor={\@author},
    	pdfsubject={\imprimirpreambulo},
	    pdfcreator={LaTeX with abnTeX2},
		pdfkeywords={abnt}{latex}{abntex}{abntex2}{plano de trabalho}, 
		colorlinks=true,       		% false: boxed links; true: colored links
    	linkcolor=blue,          	% color of internal links
    	citecolor=blue,        		% color of links to bibliography
    	filecolor=magenta,      		% color of file links
		urlcolor=blue,
		bookmarksdepth=4
}
\makeatother
% --- 

% --- 
% Espaçamentos entre linhas e parágrafos 
% --- 

% O tamanho do parágrafo é dado por:
\setlength{\parindent}{1.3cm}

% Controle do espaçamento entre um parágrafo e outro:
\setlength{\parskip}{0.2cm}  % tente também \onelineskip

\usepackage{fancyhdr}
\usepackage{lipsum} % Apenas para gerar texto fictício

% Configuração do cabeçalho especial para a primeira página
\fancypagestyle{primeirapagina}{
    \fancyhf{} % Limpa todos os campos de cabeçalho e rodapé
    \renewcommand{\headrulewidth}{0pt} % Remove a linha horizontal do topo
    
    % Define as logos alinhadas na vertical por sua base ([b]) ou centro ([c])
    \fancyhead[L]{\includegraphics[height=1.5cm, valign=b]{Images/logo-ufu.png}}
    \fancyhead[R]{\includegraphics[height=1.2cm, valign=b]{Images/logo-facom.png}}
    
    % Ajusta a distância do cabeçalho para o texto nesta página
    \setlength{\headsep}{1,5cm} 
}

% ---
% compila o indice
% ---
\makeindex
% ---

%-----------------------------------------
\begin{document}

% Ativa o cabeçalho das logos apenas nesta primeira página
\thispagestyle{primeirapagina}

% Retira espaço extra obsoleto entre as frases.
\frenchspacing 

\textbf{Disciplina:} Selecionar a opção correta conforme sua matrícula
%\textbf{Disciplina:} FACOM31701 - Trabalho de Conclusão de Curso 1 %BSI-UDI
%\textbf{Disciplina:} FACOM32703 - Projeto de Graduação I - %BSI-UDI Nova
%\textbf{Disciplina:} FACOM31703 - Trabalho de Conclusão de Curso 1 %BSI-MC (PPC-2016)
%\textbf{Disciplina:} FACOM33701 - Trabalho de Conclusão de Curso 1 %BSI-MC (PPC-2022)
%\textbf{Disciplina:} GBC072 - Projeto de Graduação 1 - %BCC
%\textbf{Disciplina:} FACOM31704 - Metodologia Científica em Computação e Informática %%BCC Nova

\textbf{Professor:} Selecionar a alternativa conforme o professor responsável pela disciplina 
%\textbf{Professor:} Prof. Dr. Rafael Dias Araújo
%\textbf{Professor:} Prof. Dr. Victor Sobreira

\textbf{Estudante:} Nome Completo

\textbf{Orientador(a):} Prof.(a) Nome Completo do Professor(a) (se, houver, indicar co-orientador abaixo)

%\textbf{Coorientador(a):} Prof.(a) Nome Completo do(a) Professor(a) Coorientador

\textbf{Tipo de trabalho: } Selecionar conforme modelo sugerido
%\textbf{Tipo de trabalho: } Projeto de Pesquisa Tradicional OU Relatório Técnico Profissional OU Desenvolvimento de Software OU Plano de Negócios

% ------------------ Título -----------------------
\section*{\centerline{\Large \textbf{Entrega 1 (E1) - Visão geral sobre o tema escolhido}}}
% -------------------------------------------------

%------------------- Corpo de dados --------------
\vspace*{0.5cm}

\textbf{Escreva um documento de até três páginas sobre o tema do TCC - contextualização, relevância, alguns exemplos de trabalhos relacionados ao tema, os objetivos gerais e uma breve descrição preliminar das ferramentas computacionais que poderão ser utilizadas durante o trabalho. O documento deverá ser organizado da seguinte forma:}

Parágrafo 1 - contextualização do tema, ou seja, uma descrição sobre o ambiente/cenário que o tema está inserido. Busque, por exemplo, por monografias defendidas por estudantes da FACOM no repositório UFU \cite{ReposiorioUFU} e outras instituições. Em geral, a contextualização do tema se encontra nos primeiros parágrafos da Seção (ou Capítulo) 1.

Parágrafo 2 - discuta a relevância do tema para a sociedade e/ou para a academia. Uma boa forma de convencer o leitor sobre a relevância do tema envolve o uso de dados, fatos ou estatísticas. Sempre utilize referências para comprovar a origem das informações.

Parágrafo 3 - liste e discuta brevemente alguns trabalhos relacionados sobre o tema. Em linhas gerais, quais problemas foram abordados por tais trabalhos? Uma estrutura muito comum para esse tipo de parágrafo é a seguinte: ``Os trabalhos X, Y e Z fornecem resultados que mostram a importância do tema. Já os trabalhos A, B e C indicam que os métodos propostos possuem bons resultados em conjuntos de dados semelhantes.'' Novamente, é importante buscar por monografias defendidas, por exemplo, por estudantes da FACOM ou outras instituições, bem como trabalhos publicados em conferências e revistas qualificadas, para observar como os demais autores redigiram este texto.

Parágrafo 4 - enuncie os objetivos gerais do trabalho. Os objetivos poderão mudar ao longo do processo. A ideia aqui é fornecer um primeiro rascunho sobre o que será feito. Lembre-se das aulas e do livro de referência \cite{metodologiaWazlawick}, especialmente da Seção 6.1 do livro em sua terceira edição. A redação do objetivo deve seguir algumas regras.

Parágrafo 5 - apresente algoritmos, softwares, hardwares e/ou métodos que possivelmente serão utilizados durante o desenvolvimento do trabalho. O conjunto de ferramentas poderá mudar ao longo do processo. A ideia aqui é fornecer um primeiro rascunho sobre as ferramentas que acreditam que serão utilizadas.

Parágrafo 6 - elabore, em conjunto com seu orientador, um parágrafo que justifique o estilo de maturidade de pesquisa que será adotado em seu TCC, conforme nomenclatura definida no Capítulo 5 do livro de referência \cite{metodologiaWazlawick}. Neste parágrafo, justifique a resposta fornecida no campo ``Tipo do Trabalho de TCC'', classificando seu trabalho como sendo Científico ou Tecnológico.

Parágrafo 7 - apresente um resumo sobre os Capítulos 3 e 4 do livro de referência \cite{metodologiaWazlawick}.

Parágrafo 8 - ao final desta entrega, declare o uso ou não de Inteligência Artificial para produção deste trabalho.

Observações:

\begin{enumerate}
    \item No caso do \textbf{relatório técnico profissional}, o texto da entrega deve estar contextualizando a empresa, o escopo do relatório (sobre o que será relatado) e a relevância do relato, conectando com os aprendizados no curso.
    \item No caso de \textbf{desenvolvimento de um software}, foque na contextualização do problema a ser resolvido, na viabilidade e demanda pela solução que será proposta.
    \item No caso de um \textbf{plano de negócios}, o texto também deve contextualizar o negócio e o(s) problema(s) que ele potencialmente soluciona e trazer uma visão geral da proposta de negócio.
    \item É de suma importância que todos os parágrafos escritos sejam coesos. Ou seja, o texto como um todo precisa fazer sentido;
    \item Use o formato de citações do Latex para apresentar as referências em todo o texto. Deve ser entendido como as citações funcionam no Latex, em particular com o BibTeX. 
    \item O posicionamento de citações no texto e orientações da ABNT serão discutidos nas aulas. Lembre-se que no modelo fornecido para a entrega 1, além do arquivo \texttt{main.tex} utilizado para redigir seu texto, há um arquivo chamado \texttt{ref.bib} com alguns exemplos de referências. Um tutorial sobre o overleaf está disponível em \cite{TutorialOverleaf}. 
    \item Junto da entrega, o estudante deverá anexar um comprovante de concordância do seu orientador quanto ao texto sendo submetido. Como alternativas preferenciais de comprovação sugerimos: anexar na entrega a cópia do e-mail do seu orientador, onde ele manifesta autorização para submeter o texto; ou uma assinatura eletrônica do orientador no PDF submetido, preferencialmente pelo sistema \href{https://www.gov.br/governodigital/pt-br/identidade/assinatura-eletronica}{GovBr} (ou outra plataforma de assinaturas conhecida e com validação pública).
\end{enumerate}


\chapter{Tema Escolhido}

Elabore aqui o texto para entrega, conforme relação de parágrafos proposta no enunciado.

Ao final desta entrega, apresentar um resumo sobre os Capítulos 3 e 4 do livro de referência \cite{metodologiaWazlawick}.

Inclua um parágrafo declarando o uso ou não de Inteligência Artificial para produção deste trabalho. A declaração deve conter a(s) ferramenta(s) utilizadas e como/onde foram utilizadas.

\bibliography{ref}

\end{document}
